\documentclass{report}

% ---------- Packages ----------
\usepackage[margin=1in]{geometry}
\usepackage{titlesec}
\usepackage{enumitem}
\usepackage{xcolor}
\usepackage{tabularx}
\usepackage{longtable}
\usepackage{booktabs}
\usepackage{hyperref}
\usepackage[T1]{fontenc}
\usepackage{parskip}
\usepackage{fancyhdr}
\usepackage{graphicx}
\usepackage{url}

% Define custom colors
\definecolor{navy}{RGB}{0, 32, 96}
\definecolor{muted}{RGB}{100, 100, 100}

% ---------- Section styling ----------
\titleformat{\section}
  {\normalfont\Large\bfseries\color{navy}}{\thesection}{0.6em}{}
\titleformat{\subsection}
  {\normalfont\large\bfseries\color{navy}}{\thesubsection}{0.6em}{}
\titlespacing*{\section}{0pt}{18pt}{8pt}
\titlespacing*{\subsection}{0pt}{12pt}{6pt}

% ---------- Header/Footer ----------
\pagestyle{fancy}
\fancyhf{}
\renewcommand{\headrulewidth}{0pt}
\fancyfoot[C]{\small\color{muted}Willow --- MVP Project Proposal}
\fancyfoot[R]{\small\color{muted}\thepage}

%% ==================================================

\begin{document}

% ===== FORMATTED TITLE PAGE =====
\begin{titlepage}
    \centering
    \vspace*{0.5cm}
    
    {\small\bfseries\color{navy} UCS503 - Software Engineering}\\[1.5cm]
    
    {\Huge\bfseries\color{navy} Willow}\\[0.5cm]
    {\large\bfseries An Image-Based Harmful-Ingredient Detection System\\for the Indian Market}\\[2cm]
    
    \begin{minipage}{0.85\textwidth}
        \centering
        \textbf{Prepared by:}\\[0.2cm]
        \texttt{[1024160001]} Aditya Gupta \\
        \texttt{[1024160019]} Akanskha Thakur\\[0.6cm]
        
        \textbf{Group:} \texttt{[3p11]} \,|\, BE Third Year, Computer Engineering\\[0.6cm]
        
        \textbf{Advisor:} Dr. Raghav B. Venkataramaiyer\\[0.6cm]
        
        \textbf{Submitted to:} Jalani Asif
    \end{minipage}
    
    \vfill
    
    {\large\bfseries Mid-Semester Evaluation}\\[0.2cm]
    {\bfseries Department of Computer Science and Engineering}\\[0.15cm]
    Thapar Institute of Engineering and Technology, Patiala
    
    \vspace*{0.5cm}
\end{titlepage}

% ===== TABLE OF CONTENTS PAGE =====
\newpage
\thispagestyle{plain}
\tableofcontents
\newpage

% ===== CHAPTER 1: INTRODUCTION =====

\chapter{Introduction}

\section{Background and Context}

Packaged and locally sold food products in the Indian market frequently
carry ingredient information that is incomplete, inconsistent, or entirely
absent from a consumer's point of view. Rising health consciousness among
Indian consumers has increased demand for tools that explain what a
product actually contains, yet the tools available today are not built
for this market. Global ingredient-scanning applications such as Yuka
\cite{yuka} rely on barcode-to-database lookups and score ingredients
against EU/US additive standards, leaving unbranded, regional, and
small-vendor Indian products -- a large share of the market -- effectively
unscanned.

This gap is not hypothetical. In April 2024, Gujarat's Food and Drug
Control Administration confiscated over 60,000 kg of adulterated spices,
including chilli powder, turmeric, and coriander powder, in a single
enforcement drive \cite{gujaratspiceseizure}. That same month, Hong Kong
and Singapore both suspended sales of spice-mix products from two major
Indian brands, MDH and Everest, after detecting the carcinogenic
pesticide ethylene oxide beyond permissible limits \cite{mdheverestban}.
Food safety and ingredient transparency are therefore live, documented
concerns in the Indian market, not a theoretical justification for this
project.

\subsection{Problem Statement}

Existing ingredient-scanning solutions available to Indian consumers are
inadequate for three concrete reasons:

\begin{enumerate}
  \item \textbf{Barcode dependency excludes local products.} Applications
    that rely on barcode-to-database lookups fail systematically for
    unbranded and regional Indian products, which are rarely
    barcode-indexed in the first place.
  \item \textbf{Many local products carry no printed ingredient list at
    all.} Local sweets, bakery items, and loose-packaged snacks are
    common in the Indian market and frequently have no ingredient panel
    to read.
  \item \textbf{Scoring standards are not localized.} Existing tools
    apply EU/US additive standards rather than FSSAI guidelines
    \cite{fssaicompendium}, and do not account for allergens or
    ingredients specific to Indian diets.
\end{enumerate}

\section{Project Objectives}

\begin{enumerate}
  \item \textbf{Objective 1:} Extract a structured ingredient list from a
    photograph of a packaged food product, without depending on a
    barcode lookup, measurably tested against a set of real product
    photographs.
  \item \textbf{Objective 2:} Flag ingredients known to be harmful and
    present their associated health effects in plain, non-technical
    language, sourced from a curated ingredient database rather than
    generated live.
  \item \textbf{Objective 3:} Handle the common case of unbranded or
    local products with no printed ingredient list by classifying the
    product into a known category and presenting a clearly labelled,
    generic risk estimate.
  \item \textbf{Objective 4:} Suggest healthier alternative products
    within the same category and a comparable price range.
  \item \textbf{Objective 5:} Build the system so that new product
    categories and data sources can be added through configuration
    rather than a redesign of the extraction or classification
    pipeline.
\end{enumerate}

These objectives were scoped deliberately to be measurable within a
one-semester, two-to-three person team timeline, rather than attempting
the full long-term vision (additional categories, crowdsourced data,
regional-language support) within the MVP itself.

% ===== CHAPTER 2: METHODOLOGY & DESIGN =====

\chapter{Methodology and Design}

\section{System Architecture}

\subsection{High-Level Design}

The system follows a modular pipeline: \emph{Photo Upload} $\rightarrow$
\emph{Extraction Service} $\rightarrow$ \emph{Classification Engine}
$\rightarrow$ \emph{Alternatives Engine} $\rightarrow$ \emph{Results}. A
separate Category Configuration store (categories, ingredients, aliases,
harm rules, and alternatives data) feeds the Classification and
Alternatives modules, so that category-specific behaviour is
data-driven rather than hard-coded into the pipeline logic. This was a
deliberate design decision: the Minimum Viable Product (MVP) is scoped to
a single category (food), but the architecture is built so that
additional categories can be added later through configuration and data
alone.

For the current phase, the backend is implemented as a modular
monolith -- clearly separated internal modules for extraction,
classification, and alternatives -- rather than deployed microservices,
to avoid infrastructure overhead that a one-semester student team cannot
absorb, while preserving a clean path to splitting into independent
services later if required.

\subsection{Technical Stack}

\begin{itemize}
  \item \textbf{Framework/Language:} Python (Flask, for the backend API
    layer) and Flutter/Dart (for the mobile frontend).
  \item \textbf{Database:} PostgreSQL is planned for the category
    configuration store (categories, ingredients, aliases, alternatives);
    the current extraction-only prototype does not yet require a
    persistent database.
  \item \textbf{Tools:} Git for version control, VS Code for development,
    Overleaf/\LaTeX{} for documentation, and draw.io for architecture,
    use-case, ER, and activity diagrams.
  \item \textbf{Libraries:} \texttt{google-genai} and \texttt{anthropic}
    (vision-capable extraction APIs), \texttt{pytesseract} and
    \texttt{Pillow} (local OCR fallback), \texttt{flask-cors} (API
    layer), and \texttt{http}/\texttt{image\_picker} (Flutter HTTP
    client and camera capture).
\end{itemize}

\section{Implementation Details}

\subsection{Module 1: Ingredient Extraction Pipeline}

The extraction layer is built behind a single abstract interface
(\texttt{IngredientExtractor}), so that the underlying vision provider
can be swapped without any change to downstream code. Three concrete
implementations exist at this stage: a local, free Tesseract-based
extractor (used for offline development and testing without API cost), a
Gemini vision-based extractor, and a Claude vision-based extractor,
allowing direct accuracy and latency comparison between providers on the
same input.

A deliberate architectural boundary was drawn early in this module: the
vision API's only responsibility is reading text off the image. It does
not judge, classify, or explain any ingredient. This decision was made
because large language models are prone to subtle inaccuracies on
exactly the kind of specific, falsifiable health claims this application
makes, and an incorrect claim shown to a real user is a liability
concern, not merely a UX one. Classification against a curated,
sourced database is therefore treated as a separate module (planned for
the next phase), so that every harmful-ingredient flag shown to the user
can be traced back to a specific, checkable source rather than to a
live model inference.

A parsing layer sits between raw model output and the rest of the
system: it splits comma/semicolon-separated ingredient text while
correctly preserving nested parenthetical groups (for example,
\texttt{"Raising Agents (INS 500(ii))"} is treated as a single
ingredient, not split at the inner parenthesis), strips common
label prefixes such as \texttt{"INGREDIENTS:"}, and normalises
whitespace. This parser is covered by a unit test suite (Section 3.1)
and was verified against a realistic synthetic Indian product label
containing INS-coded additives before being used on real photographs.

\subsection{Module 2: Flutter Frontend Integration}

To connect the existing Flutter frontend to the extraction pipeline for
the current evaluation milestone, a minimal REST API was built around
the extraction module using Flask. A single endpoint,
\texttt{POST /extract}, accepts a multipart image upload and an optional
provider selector, and returns a JSON response containing the parsed
ingredient list, the raw extracted text, a confidence score, and the
provider used. Error conditions (unsupported file type, missing image
field, or a provider-level extraction failure) return a 400 response
with a specific, human-readable error message rather than a generic
failure.

On the Flutter side, the \texttt{image\_picker} package is used to
capture a photograph directly from the device camera, and the
\texttt{http} package's \texttt{MultipartRequest} is used to upload it
to the backend. A practical integration detail worth documenting: the
correct backend base URL differs by run target (\texttt{10.0.2.2} for
the Android emulator, \texttt{127.0.0.1} for the iOS simulator, and the
development machine's LAN IP address for a physical device on the same
network) -- a common point of confusion when connecting a mobile client
to a locally running backend during development.

% ===== CHAPTER 3: RESULTS & EVALUATION =====

\chapter{Results and Evaluation}

\section{Testing Methodology}

\begin{itemize}
  \item \textbf{Unit Testing:} The ingredient-list parser is covered by a
    dedicated unit test suite exercising comma/semicolon splitting,
    nested parenthetical groups, label-prefix stripping, whitespace
    normalisation, and empty-input handling.
  \item \textbf{Integration Testing:} The full extraction pipeline
    (photo $\rightarrow$ extraction $\rightarrow$ parsing) was verified
    end-to-end against a synthetic product label, and the Flask API
    layer was verified against the same pipeline using real HTTP
    requests, including both successful extractions and the two
    documented error paths.
  \item \textbf{Performance Testing:} A comparison harness was built to
    run every configured extraction provider (Tesseract, Gemini, Claude)
    against a shared set of product photographs and report per-provider
    latency, and, where a manually transcribed ground truth is supplied,
    precision and recall of the extracted ingredient list.
  \item \textbf{User Testing:} Not yet conducted. A pilot with 10-15
    testers, comparing system output against manual expert review of the
    same product photographs, is planned once the classification engine
    (Chapter 4) is complete.
\end{itemize}

\section{Performance Metrics}

The results below are preliminary smoke-test results from the current
extraction-only milestone, not the full pilot evaluation, which depends
on the classification engine described in Section 4.3.

\begin{table}[h]
\centering
\begin{tabular}{|l|r|}
  \hline
  \textbf{Metric} & \textbf{Result} \\
  \hline
  Parser unit tests passing & $10/10$ \\
  \hline
  Ingredients correctly parsed (synthetic label test) & $8/8$ \\
  \hline
  Precision (synthetic label, deliberate-mismatch check) & $1.00$ \\
  \hline
  Recall (synthetic label, deliberate-mismatch check) & $0.89$ \\
  \hline
  API error-handling paths verified & $2/2$ \\
  \hline
  Real-world accuracy (10-20 photo test set, mixed conditions) & Pending \\
  \hline
\end{tabular}
\caption{Preliminary Extraction Pipeline Results}
\label{tab:performance}
\end{table}

\section{Analysis}

The objectives met so far are Objective 1 (extraction pipeline, built
and unit-tested) and the frontend integration required to demonstrate it
end-to-end from the Flutter application. Objectives 2 through 4
(harmful-ingredient flagging, the unbranded-product fallback, and
alternative suggestions) depend on the classification engine and
category-configuration database, which are the next planned milestone
and have not yet been implemented.

Two challenges were encountered and resolved during this phase. First,
an early design question was whether the vision API used for extraction
should also judge which ingredients are harmful. This was rejected in
favour of a strict separation (Section 2.2.1), since a live model
judgment cannot be reliably traced to a specific, checkable source for a
health-related claim. Second, connecting the Flutter frontend to a
locally running backend surfaced the common but under-documented issue
of platform-specific base URLs; this was resolved by explicitly
documenting the correct host address for each run target rather than
assuming a single URL would work across emulator, simulator, and
physical-device testing.

The remaining, not-yet-resolved challenge is real-world extraction
accuracy: the pipeline has only been verified against a single synthetic
label so far. Testing against a genuinely varied set of real product
photographs (clear, blurry, curved packaging, and low light) is the
immediate next step and is required before any accuracy claim can be
made with confidence.

% ===== CHAPTER 4: CONCLUSION =====

\chapter{Conclusion}

\section{Summary of Work}

At this mid-semester stage, the project has completed requirements
scoping, system architecture design, and a working, tested ingredient
extraction pipeline with two interchangeable vision-API providers and a
local free-tier fallback, connected end-to-end to the existing Flutter
frontend through a minimal REST API. Objective 1 is functionally
complete pending real-world accuracy testing; Objectives 2 through 5 are
designed but not yet implemented.

\section{Key Findings}

\begin{itemize}
  \item \textbf{Finding 1:} An image-first extraction approach, rather
    than a barcode-lookup approach, is architecturally necessary to
    cover unbranded and local Indian products, which make up a
    significant share of the market existing tools do not reach.
  \item \textbf{Finding 2:} Separating ingredient extraction (a reading
    task, delegated to a vision API) from harmful-ingredient
    classification (a judgment task, delegated to a curated, sourced
    database) removes a specific and serious hallucination risk that a
    single-model approach would carry on health-related claims.
  \item \textbf{Finding 3:} Practical mobile-to-backend integration
    details, such as platform-specific base URLs, are a recurring but
    often undocumented source of friction when connecting a Flutter
    client to a locally hosted development server, and are worth
    documenting explicitly for the rest of the team.
\end{itemize}

\section{Future Work and Recommendations}

\begin{enumerate}
  \item Build the classification engine and the category-configuration
    database (categories, curated ingredients, aliases, and harm rules),
    seeded from Open Food Facts \cite{openfoodfacts} and cross-referenced
    against the FSSAI Compendium on Food Safety and Standards (Food
    Additives) Regulation \cite{fssaicompendium}.
  \item Implement the category-based fallback for unbranded products
    with no printed ingredient list, including the generic risk-profile
    disclaimer described in the project's requirements documentation.
  \item Build the healthier-alternatives engine and the full results
    screen in the Flutter application.
  \item Conduct the planned pilot test (10-15 users) and collect the
    evaluation metrics defined for this project: extraction accuracy,
    flagging precision, and time-to-result.
  \item Stretch goals, to be attempted only if the core scope above is
    complete with time remaining: a second product category (cosmetics
    and personal care), regional-language OCR support, and integration
    with the FSSAI FoSCoS license lookup as an additional trust signal
    for unbranded products.
\end{enumerate}

\section{Final Remarks}

The mid-semester phase of this project has validated the core technical
approach -- image-based extraction with a provider-agnostic interface --
against real code and real (if preliminary) test results, rather than
design alone. The most significant remaining risk is not the extraction
pipeline itself, which is functioning, but the accuracy of extraction
across the genuinely messy conditions of real Indian product packaging,
which has not yet been measured. Closing that gap, and building the
classification engine that turns an extracted ingredient list into an
actionable, trustworthy health signal for the user, are the primary
goals for the remainder of the semester.

% ===== BIBLIOGRAPHY =====

\begin{thebibliography}{99}
\addcontentsline{toc}{chapter}{Bibliography}

\bibitem{yuka}
Yuka. \textit{Yuka}. \url{https://yuka.io}. Accessed 16 September 2026.

\bibitem{fssaicompendium}
Food Safety and Standards Authority of India. \textit{Compendium on Food
Safety and Standards (Food Products Standards and Food Additives)
Regulation}. 2011. \url{https://fssai.gov.in/cms/Compendium-FSS-FPS-FA.php}.
Accessed 16 September 2026.

\bibitem{openfoodfacts}
Open Food Facts. \textit{Open Food Facts --- Data, API and SDKs}.
\url{https://world.openfoodfacts.org/data}. Accessed 16 September 2026.

\bibitem{gujaratspiceseizure}
Daijiworld. \textit{Gujarat: Food Safety Raid Uncovers 1,500 kg of
Adulterated Paneer in Ahmedabad}. 2025.
\url{https://www.daijiworld.com/index.php/news/newsDisplay?newsID=1268855}.
Article reports the April 2024 Gujarat FDCA seizure of over 60,000 kg of
adulterated spices. Accessed 16 September 2026.

\bibitem{mdheverestban}
Deccan Herald. \textit{After Singapore, Hong Kong Bans Everest Fish Curry
Masala, MDH Spices}. 2024.
\url{https://www.deccanherald.com/business/companies/after-singapore-hong-kong-bans-everest-fish-curry-masala-mdh-spices-2988741}.
Accessed 16 September 2026.

\end{thebibliography}

\end{document}